\section{Arquitectura general}
\label{sec:arquitectura_general}

La plataforma está organizada en tres componentes principales que se
reutilizan parcialmente entre los dos modos de despliegue:

\begin{enumerate}
  \item \textbf{Aplicación móvil}: punto de entrada para el operario en
  el escenario en la nube (captura de foto, autenticación, historial de
  resultados).
  \item \textbf{CLOUD\_MODE}: conjunto de microservicios desplegados en
  Railway (API Gateway, servicio de autenticación, servicio de
  clasificación, broker MQTT) orientado a múltiples usuarios.
  \item \textbf{LOCAL\_MODE}: un servidor FastAPI de un solo proceso,
  sin autenticación ni TLS, pensado para correr en la misma red cerrada
  que el ESP32 y la cámara, sin salida a internet.
\end{enumerate}

Los dos modos no interoperan entre sí: un despliegue dado opera en uno
u otro, según el entorno disponible (infraestructura administrada vs.
banco de pruebas aislado). La Tabla \ref{tab:comparacion} resume las
diferencias principales.

\begin{table}[H]
\centering
\caption{Comparación entre CLOUD\_MODE y LOCAL\_MODE}
\label{tab:comparacion}
\begin{tabular}{@{}p{2.1cm}p{2.7cm}p{2.7cm}@{}}
\toprule
 & \textbf{CLOUD\_MODE} & \textbf{LOCAL\_MODE} \\
\midrule
Autenticación & JWT (RS256) & Ninguna \\
Transporte & HTTPS + MQTT/TLS & HTTP plano \\
Usuarios & Múltiples, con cuentas & Un solo operario \\
Almacenamiento & Objeto remoto (S3) & En memoria, un proceso \\
Alcance de red & Internet & LAN cerrada (hotspot) \\
Hardware físico & ESP32-S3, solo muestra resultado & ESP32-WROOM-32, controla banda y huella \\
\bottomrule
\end{tabular}
\end{table}
